\subsection{Ciclos primos asociados sobre los preesquemas algebraicos.}
\label{subsection:IV.2.4.2-fr}

\begin{proposition}[4.2.1]
\EGAReaderTarget{ega4.c04.prop.4-2-1}{IV.2.4.2.1-fr}
Sean $K$ y $L$ dos extensiones de un cuerpo $k$, tales que $K\otimes_k L$
sea noetheriano. Entonces los ideales primos asociados a $K\otimes_k L$ son
mínimos y, si $E$ es el cuerpo residual del anillo local de uno de esos ideales,
se tiene
\begin{equation}
\tag{4.2.1.1}
\label{IV.2.4.2.1.1-fr}
\operatorname{deg.tr}_K E=\operatorname{deg.tr}_k L,
\qquad
\operatorname{deg.tr}_L E=\operatorname{deg.tr}_k K
\end{equation}
de donde
\begin{equation}
\tag{4.2.1.2}
\label{IV.2.4.2.1.2-fr}
\operatorname{deg.tr}_k E
=\operatorname{deg.tr}_k K+\operatorname{deg.tr}_k L.
\end{equation}
\end{proposition}

\begin{proof}
Se sabe que $K$ es una extensión algebraica de una extensión puramente
trascendente $K'=k(\mathbf t)$, donde $\mathbf t=(t_\alpha)_{\alpha\in I}$ es
una familia de indeterminadas; el anillo $k[\mathbf t]\otimes_k L=L[\mathbf t]$
es íntegro, por tanto lo mismo ocurre con $K'\otimes_k L$, que es un anillo de
fracciones de este, y el cuerpo de fracciones de $K'\otimes_k L$ es
$L(\mathbf t)$; se tiene el diagrama conmutativo de homomorfismos canónicos

\oldpage[IV-2]{55}

\begin{equation}
\tag{4.2.1.3}
\label{IV.2.4.2.1.3-fr}
\xymatrix@C=3.5em@R=2.8em{
K \ar[r] & K\otimes_k L \ar[r] & K\otimes_{K'}L(\mathbf t)\\
K' \ar[u] \ar[r] & K'\otimes_k L \ar[u] \ar[r] & L(\mathbf t) \ar[u]\\
k \ar[u] \ar[r] & L \ar[u] &
}
\end{equation}

Como $K$ es fielmente plano sobre $K'$, $K\otimes_k L
=K\otimes_{K'}(K'\otimes_k L)$ es fielmente plano sobre $K'\otimes_k L$,
por tanto $K'\otimes_k L$ es noetheriano $(\mathbf{0}_{\mathrm I},6.5.2)$; además
$K'\otimes_k L$ se identifica con un subanillo de $K\otimes_k L$; la traza
sobre $K'\otimes_k L$ de un ideal primo $\mathfrak p$ asociado a $K\otimes_k L$
es el ideal $0$, pues un elemento $\neq0$ de $K'\otimes_k L$ no es divisor de
cero en $K\otimes_k L$ $(\mathbf{0}_{\mathrm I},6.3.4$ y Bourbaki,
\emph{Alg. comm.}, cap.~IV, \S~1, n\textsuperscript{o}~1, cor.~3 de la
prop.~2). Como además $K$ es algebraico sobre $K'$, $K\otimes_k L$ es un
álgebra entera sobre $K'\otimes_k L$, y los ideales primos de $K\otimes_k L$
que inducen $0$ en $K'\otimes_k L$ carecen necesariamente de relación de
inclusión mutua (Bourbaki, \emph{Alg. comm.}, cap.~V, \S~2,
n\textsuperscript{o}~1, cor.~1 de la prop.~1); esto demuestra la primera
afirmación (Bourbaki, \emph{Alg. comm.}, cap.~IV, \S~1,
n\textsuperscript{o}~3, cor.~1 de la prop.~7). Además, el cuerpo residual $E$
de $\mathfrak p$ es algebraico sobre el cuerpo residual del ideal $(0)$ de
$K'\otimes_k L$, es decir, $L(\mathbf t)$; por tanto
$\operatorname{deg.tr}_k E=\operatorname{deg.tr}_k L(\mathbf t)
=\operatorname{Card}(I)=\operatorname{deg.tr}_k K$, en otras palabras, se
tiene la primera relación (4.2.1.1); intercambiando los papeles de $K$ y $L$,
se obtiene la segunda relación (4.2.1.1), de donde se sigue (4.2.1.2).
\end{proof}

\begin{corollary}[4.2.2]
\label{IV.2.4.2.2-fr}
Bajo las hipótesis de \textup{(3.3.6)}, si los preesquemas $T_{x,y}$ son localmente
noetherianos, no tienen ciclos primos asociados inmersos.
\end{corollary}

\begin{corollary}[4.2.3]
\label{IV.2.4.2.3-fr}
Bajo las hipótesis de \textup{(3.3.6) (resp. (3.3.7))}, si los $T_{x,y}$ son
localmente noetherianos y si $\mathcal F$ y los $\mathcal G_s$ \textup{(}para $s\in S$\textup{)}
\textup{(resp.} $\mathcal F$ y $\mathcal G$\textup{)} carecen de ciclos primos asociados
inmersos, lo mismo ocurre con $\mathcal F\otimes_S\mathcal G$.
\end{corollary}

Esto se deduce de (4.2.2), (3.3.2) y de la demostración de (3.3.6). En
particular, como todo preesquema sobre un cuerpo $k$ es \emph{plano} sobre $k$,
se ha demostrado así la afirmación (i) de la
